\chapter{Introduction}
\label{sec:introduction}
Many fundamental questions in social sciences and epidemiology are 
causal in nature focusing on the \emph{effect} of an \emph{intervention} on an outcome. 
Examples include whether increasing the minimum wage affects employment \cite{cardMinimumWagesEmployment1993} 
or whether initiating hormone therapy changes the risk of coronary heart disease in postmenopausal 
women \cite{hernanObservationalStudiesAnalyzed2008}. In order to answer such questions, it is not sufficient to study 
associations alone. Rather, causal inference requires assumptions about the underlying data-generating process (DGP)
and a clear definition of the causal quantity of interest \cite{greenlandCausalDiagramsEpidemiologic1999, hernanUsingBigData2016}. 
While randomized controlled trials are often regarded as the gold standard for estimating causal effects, 
they are usually infeasible, unethical, or expensive in practice. As a result, researchers often have to rely on 
observational data to estimate causal effects, which introduces additional challenges due to potential confounding.

This thesis focuses on a specific observational longitudinal study design with 
discrete time points \cite{gillCausalInferenceComplex2001}, which involves
data with repeated covariate measurements of the same units, a time-varying binary treatment, 
a terminal outcome at the end of the observation period and a lack of anticipation effects, meaning 
that treatments at future time points do not impact potential outcomes at the current time point \cite{arkhangelskyCausalModelsLongitudinal2024}. 
Instances of such designs can be found in psychology \cite{molinaAssociationStimulantTreatment2023}, sociology \cite{bacakMarginalStructuralModels2015}, and epidemiology \cite{bodnarMarginalStructuralModels2004}. 
Moreover, the growing availability of intensive longitudinal data through techniques like experience sampling, 
mobile evaluations, and wearable devices has made methods for longitudinal causal 
inference increasingly important \cite{myin-germeysExperienceSamplingMethodology2018, shiMetalearningMethodEstimation2025a}.

Estimating causal effects in longitudinal settings with time-varying treatments and confounders is 
particularly challenging due to the presence of treatment-confounder feedback, where treatments at 
previous time points can affect confounders at subsequent time points, which in turn can affect future 
treatments and the outcome. In the presence of such feedback loops, naive standard regression methods that 
adjust for 
post-treatment confounders at each time point can produce biased estimates of causal effects. %\cite{hernanCausalInferenceWhat2023booksection}
To address this issue, a family of methods known as the g-methods have been developed \cite{robinsEstimationCausalEffects2008, naimiIntroductionMethods2017}. 
However, these methods rely on the assumption of no unmeasured confounding or sequential ignorability. This assumption states 
that all confounders of the treatment-outcome relationship are measured and included in the analysis,
which is often violated in observational studies.

There are two types of unmeasured confounding that can occur in longitudinal 
studies: time-varying and time-invariant. In this thesis, the focus is on the latter, 
which are also referred to as missing baseline confounders. These are 
unmeasured covariates that affect the treatments and outcome at all time 
points but do not change over time. Examples of such baseline covariates may 
include genetic predisposition, personality traits, or other stable 
characteristics of individuals during the span of the study. 

Two methods have been proposed in the recent literature to estimate the 
causal effects of time-varying treatments in the presence of time-invariant 
unmeasured confounders. The first is the fixed-effect marginal structural model (FE-MSM) 
proposed by \citeA{blackwellEffectPoliticalAdvertising2024}, 
which combines the marginal structural model (MSM) 
\cite{robinsMarginalStructuralModels2000} with unit fixed effects. The second method 
is the sequential Gaussian process latent variable model (SeqGPLVM) introduced by 
\citeA{hattSequentialDeconfoundingCausal2024}. This method may be interpreted as the 
sequential extension of the deconfounder theory \cite{wangBlessingsMultipleCauses2019}.

%where 
%a latent variable is used to estimate a statistically sufficient substitute for the unmeasured confounders. This allows the method to retain the sequential ignorability 
%assumption to estimate any causal target of interest in a downstream analysis.

These two methods are comparable as they can be used in similar settings such as those described above. They both involve a 
unit-specific constant unobserved heterogeneity component that is supposed to 
capture the effect of the unmeasured confounders. They diverge, however, 
in their characterization of the unobserved heterogeneity component. 
Whereas FE-MSM---rooted in econometrics literature---conceptualizes this component 
as a frequentist fixed effect and a nuisance parameter which is tied to inverse 
probability of treatment weighting (IPTW) of marginal structural models, 
the machine-learning-based SeqGPLVM views it as a 
latent variable and employs a much more flexible and non-parametric approach to 
estimate that latent variable which then can be employed in a wide range of 
causal inference estimators. As a result, they make different assumptions 
about the unobserved heterogeneity component, the range of causal targets they can 
identify and the scope of estimators they can be used with. Despite 
these similarities and differences, these two methods have not been 
compared in the same paper using a common notation. This makes it 
difficult to understand where the two methods overlap in their assumptions and 
when one method might be preferred over the other. 

To fill this gap, the current thesis provides a unified presentation of these two 
methods, highlighting their assumptions with a common notation. The performance of 
the proposed methods was evaluated through a simulation study using 
the data-generating process suggested in \citeA{blackwellEffectPoliticalAdvertising2024}. 
Synthetic datasets were generated for three different ratios of individuals to time 
points under two different confounding strengths. The estimates obtained from 
both methods were compared to the true causal effects to assess their bias, 
variance and confidence interval coverage. 

Beyond the simulation study, a detailed derivation of the loss function for 
SeqGPLVM is provided, which was not fully elaborated in the original paper. 
Moreover, a modular implementation of SeqGPLVM is 
developed using the \texttt{GPyTorch} library \cite{gardnerGPyTorchBlackboxMatrixMatrix2021}. 
This implementation is designed to accommodate different treatment types and is intended for 
practical reuse in future research\footnote{Code available at \url{https://github.com/Ali-Setareh/SeqGPLVM_FEMSM}}.

The remainder of the thesis is structured as follows. Chapter ~\ref{chap:literature_review} reviews the 
relevant literature on longitudinal causal inference and the methods that exist to address confounding. Next in
Chapter ~\ref{chap:method}, the FE-MSM and the 
SeqGPLVM are described in detail, including their assumptions and estimation procedures. 
Chapter ~\ref{chap:simulation_study} presents a simulation study comparing the performance of these 
two methods under different scenarios. Finally, chapter ~\ref{chap:discussion} discusses the 
implications of the findings and suggests directions for future research.

